% ═══════════════════════════════════════════════════════════════
% ML-02c: Possessivbegleiter (MUSTERLÖSUNG)
% Spanisch Klasse 7 - Sequenz 2: Mi familia
% ═══════════════════════════════════════════════════════════════
% Identisches Layout wie AB-02c, Lösungen in ROT
% ═══════════════════════════════════════════════════════════════

\documentclass[11pt, a4paper]{article}

% ═══════════════════════════════════════════════════════════════
% SW-MODUS TOGGLE
% ═══════════════════════════════════════════════════════════════
\newif\ifswmodus
\swmodusfalse

% ═══════════════════════════════════════════════════════════════
% PAKETE
% ═══════════════════════════════════════════════════════════════

\usepackage[utf8]{inputenc}
\usepackage[T1]{fontenc}
\usepackage[ngerman]{babel}
\usepackage{helvet}
\renewcommand{\familydefault}{\sfdefault}

\usepackage[margin=2cm]{geometry}
\usepackage{enumitem}
\usepackage{tabularx}
\usepackage{booktabs}
\usepackage{multirow}

\usepackage{xcolor}
\usepackage{framed}
\usepackage{tcolorbox}
\tcbuselibrary{skins}

\usepackage{fancyhdr}
\usepackage{lastpage}
\usepackage{needspace}

% ═══════════════════════════════════════════════════════════════
% FARBEN
% ═══════════════════════════════════════════════════════════════

\definecolor{spanisch-color}{HTML}{c41e3a}
\definecolor{grau-color}{HTML}{888888}
\definecolor{hellgrau-color}{HTML}{f5f5f5}
\definecolor{orange-color}{HTML}{b35c00}
\definecolor{loesung-color}{HTML}{c62828}

\definecolor{sw-dark}{HTML}{000000}
\definecolor{sw-gray}{HTML}{666666}
\definecolor{sw-white}{HTML}{ffffff}

\ifswmodus
    \colorlet{spanisch}{sw-dark}
    \colorlet{grau}{sw-gray}
    \colorlet{hellgrau}{sw-white}
    \colorlet{orange}{sw-dark}
    \colorlet{loesung}{sw-dark}
\else
    \colorlet{spanisch}{spanisch-color}
    \colorlet{grau}{grau-color}
    \colorlet{hellgrau}{hellgrau-color}
    \colorlet{orange}{orange-color}
    \colorlet{loesung}{loesung-color}
\fi

\newcommand{\fachfarbe}{spanisch}

% ═══════════════════════════════════════════════════════════════
% VARIABLEN
% ═══════════════════════════════════════════════════════════════

\newcommand{\thema}{Possessivbegleiter}
\newcommand{\sequenz}{02}
\newcommand{\block}{c}
\newcommand{\fach}{Spanisch}
\newcommand{\klasse}{7}
\newcommand{\maxpunkte}{20}
\newcommand{\datum}{\today}

% ═══════════════════════════════════════════════════════════════
% KOPF- UND FUSSZEILE
% ═══════════════════════════════════════════════════════════════

\pagestyle{fancy}
\fancyhf{}
\renewcommand{\headrulewidth}{0.4pt}
\renewcommand{\footrulewidth}{0.4pt}

\lhead{\fach{} -- Klasse \klasse}
\chead{\textcolor{loesung}{\textbf{ML-\sequenz\block{} (MUSTERLÖSUNG)}}}
\rhead{\datum}

\lfoot{}
\cfoot{}
\rfoot{Seite \thepage\ von \pageref{LastPage}}

% ═══════════════════════════════════════════════════════════════
% BEFEHLE
% ═══════════════════════════════════════════════════════════════

\newcommand{\punktebox}[1]{\hfill\fbox{\small \_/#1}}
\newcommand{\luecke}[1]{\rule{#1}{0.4pt}}
\newcommand{\lsg}[1]{\textcolor{loesung}{\textbf{#1}}}

\newtcolorbox{merkkasten}[1][]{
    colback=hellgrau,
    colframe=\fachfarbe,
    colbacktitle=white,
    coltitle=\fachfarbe,
    boxrule=1pt,
    arc=0pt,
    left=8pt, right=8pt, top=8pt, bottom=8pt,
    fonttitle=\bfseries,
    attach boxed title to top left={yshift=-2mm, xshift=5mm},
    boxed title style={boxrule=0pt, colback=white},
    title=#1
}

\newtcolorbox{vokabelkasten}{
    colback=white,
    colframe=\fachfarbe,
    boxrule=0.5pt,
    arc=0pt,
    left=5pt, right=5pt, top=5pt, bottom=5pt
}

\newtcolorbox{hinweis}{
    colback=white,
    colframe=orange,
    colbacktitle=white,
    coltitle=orange,
    boxrule=1pt,
    arc=0pt,
    left=5pt, right=5pt, top=5pt, bottom=5pt,
    fonttitle=\bfseries,
    attach boxed title to top left={yshift=-2mm, xshift=5mm},
    boxed title style={boxrule=0pt, colback=white},
    title=Hinweis
}

\newenvironment{aufgabe}[1]{%
    \needspace{5\baselineskip}%
    \nopagebreak[4]%
    \vspace{10pt}
    \noindent\textbf{\textcolor{\fachfarbe}{Aufgabe #1}}
    \nopagebreak[4]%
    \vspace{5pt}
    \nopagebreak[4]%
    \begin{list}{}{\leftmargin=0pt}
    \item[]
}{%
    \end{list}
}

\newcommand{\es}[1]{\textcolor{\fachfarbe}{\textbf{#1}}}

% ═══════════════════════════════════════════════════════════════
% DOKUMENT
% ═══════════════════════════════════════════════════════════════

\begin{document}

% ═══════════════════════════════════════════════════════════════
% KOPFBEREICH
% ═══════════════════════════════════════════════════════════════

\begin{center}
    {\Large\bfseries\textcolor{\fachfarbe}{Arbeitsblatt: \thema}}\\[0.3em]
    {\large\textcolor{loesung}{\textbf{--- MUSTERLÖSUNG ---}}}
\end{center}

\vspace{5pt}

\noindent
\begin{tabularx}{\textwidth}{|X|X|X|}
\hline
\textbf{Name:} \lsg{Musterschüler/in} &
\textbf{Klasse:} \klasse &
\textbf{Datum:} \datum \\
\hline
\end{tabularx}

\vspace{15pt}

% ═══════════════════════════════════════════════════════════════
% MERKKASTEN: Die Possessivbegleiter
% ═══════════════════════════════════════════════════════════════

\begin{merkkasten}[Merkkasten: Die Possessivbegleiter]
Die Possessivbegleiter zeigen, wem etwas gehört:

\vspace{0.5em}

\begin{center}
\begin{tabular}{|l|c|c|c|}
\hline
\textbf{Person} & \textbf{Deutsch} & \textbf{Singular} & \textbf{Plural} \\
\hline
yo (ich) & mein/e & \es{mi} \lsg{(mein/e)} & \es{mis} \lsg{(meine)} \\
\hline
tú (du) & dein/e & \es{tu} \lsg{(dein/e)} & \es{tus} \lsg{(deine)} \\
\hline
él/ella (er/sie) & sein/e, ihr/e & \es{su} \lsg{(sein/e, ihr/e)} & \es{sus} \lsg{(seine, ihre)} \\
\hline
\end{tabular}
\end{center}

\vspace{0.5em}

\textbf{Wichtig:} Der Possessivbegleiter richtet sich nach dem \textbf{Bezugswort}!

\begin{itemize}[leftmargin=2em, itemsep=0pt]
    \item Singular: \es{mi} madre, \es{tu} padre, \es{su} hermano
    \item Plural: \es{mis} hermanos, \es{tus} abuelos, \es{sus} primos
\end{itemize}
\end{merkkasten}

\vspace{10pt}

% ═══════════════════════════════════════════════════════════════
% AUFGABE 1: mi/tu/su einsetzen
% ═══════════════════════════════════════════════════════════════

\begin{aufgabe}{1 -- Setze mi, tu oder su ein}
Ergänze den passenden Possessivbegleiter. \punktebox{4}
\end{aufgabe}

\begin{vokabelkasten}
\textit{Wähle:} \es{mi} | \es{tu} | \es{su}
\end{vokabelkasten}

\vspace{0.5em}

\noindent
a) \lsg{Mi} madre se llama Carmen. (meine Mutter)

\vspace{0.8em}

\noindent
b) ¿Cómo se llama \lsg{tu} padre? (dein Vater)

\vspace{0.8em}

\noindent
c) \lsg{Su} hermano tiene 15 años. (sein Bruder)

\vspace{0.8em}

\noindent
d) ¿Cuántos años tiene \lsg{tu} abuela? (deine Großmutter)

\vspace{15pt}

% ═══════════════════════════════════════════════════════════════
% AUFGABE 2: Singular $\rightarrow$ Plural
% ═══════════════════════════════════════════════════════════════

\begin{aufgabe}{2 -- Bilde den Plural}
Schreibe die Pluralform. \punktebox{4}
\end{aufgabe}

\begin{hinweis}
Singular: \es{mi} / \es{tu} / \es{su} $\rightarrow$ Plural: \es{mis} / \es{tus} / \es{sus}
\end{hinweis}

\vspace{0.5em}

\noindent
\textbf{Beispiel:} \es{mi hermano} $\rightarrow$ \es{mis hermanos}

\vspace{1em}

\noindent
a) \es{mi hermano} $\rightarrow$ \lsg{mis hermanos}

\vspace{0.8em}

\noindent
b) \es{tu abuelo} $\rightarrow$ \lsg{tus abuelos}

\vspace{0.8em}

\noindent
c) \es{su prima} $\rightarrow$ \lsg{sus primas}

\vspace{0.8em}

\noindent
d) \es{mi tío} $\rightarrow$ \lsg{mis tíos}

\vspace{15pt}

% ═══════════════════════════════════════════════════════════════
% AUFGABE 3: Familie beschreiben
% ═══════════════════════════════════════════════════════════════

\begin{aufgabe}{3 -- Beschreibe deine Familie}
Schreibe drei Sätze über deine Familie. Verwende \es{mi} oder \es{mis}. \punktebox{6}
\end{aufgabe}

\begin{vokabelkasten}
\textit{Redemittel:} \es{Mi madre se llama...} | \es{Mi padre tiene... años.} | \es{Mis hermanos se llaman...}
\end{vokabelkasten}

\vspace{0.8em}

\textbf{Beispiellösungen:}

\noindent
1. \lsg{Mi madre se llama Carmen. Tiene 42 años.}

\vspace{1em}

\noindent
2. \lsg{Mi padre se llama Pedro. Es profesor.}

\vspace{1em}

\noindent
3. \lsg{Mis hermanos se llaman Luis y Ana.}

\vspace{0.5em}
\textit{\textcolor{grau}{Bewertung: Je 2P pro korrektem Satz (Possessiv + Verb + Info)}}

\vspace{15pt}

% ═══════════════════════════════════════════════════════════════
% AUFGABE 4: Dialog
% ═══════════════════════════════════════════════════════════════

\begin{aufgabe}{4 -- Vervollständige den Dialog}
A fragt nach der Familie von B. Verwende \es{mi} / \es{tu} / \es{su}. \punktebox{6}
\end{aufgabe}

\begin{hinweis}
Verwende: \es{¿Cómo se llama tu...?} und \es{Mi ... se llama...}
\end{hinweis}

\vspace{0.5em}

\noindent
\textbf{A:} ¿Cómo se llama \lsg{tu} hermano?

\vspace{0.8em}

\noindent
\textbf{B:} \lsg{Mi} hermano se llama \lsg{Pablo}.

\vspace{0.8em}

\noindent
\textbf{A:} ¿Cuántos años tiene?

\vspace{0.8em}

\noindent
\textbf{B:} Tiene \lsg{doce} años.

\vspace{0.8em}

\noindent
\textbf{A:} ¿Y \lsg{tu} madre? ¿Cómo se llama?

\vspace{0.8em}

\noindent
\textbf{B:} \lsg{Mi} madre se llama \lsg{Ana}.

\vspace{0.5em}
\textit{\textcolor{grau}{Bewertung: Fragen korrekt (2P), Antworten mit Possessiv (2P), Namen/Alter korrekt (2P)}}

% ═══════════════════════════════════════════════════════════════
% PUNKTEÜBERSICHT
% ═══════════════════════════════════════════════════════════════

\vspace{20pt}
\hrule
\vspace{10pt}

\begin{flushright}
\textbf{Punkte:} \lsg{\maxpunkte} / \maxpunkte
\end{flushright}

\end{document}
